% GENERATED FILE. Edit canonical lessons, then run npm run generate:books.
% canonical-source-hash: fnv1a64:4d82424f6d53d031
% canonical-lessons: TA-C452-miti-pedal, TA-C452-carakku, TA-C452-piratecam, TA-C452-pakuti, TA-C452-urul, TA-R452-review-446-452, TA-R452-review-final-452

\chapter{Heading Home}
\label{ch:ta-heading-home}
\hlchaptermodality{452}

\begin{chapteropening}
\textbf{By the end of this chapter:} \emph{I can talk about pedalling, goods, regions and areas, and wheels rolling.}

Complete the last lesson of chapter 452: I can talk about pedalling, goods, regions and areas, and wheels rolling.
\end{chapteropening}

\section[\texorpdfstring{\ta{மிதி} (miti)}{மிதி (miti)}]{\ta{மிதி} (miti) — to pedal; to tread on}

\label{lesson:TA-C452-miti-pedal}



\begin{quote}
Before the new one: say the Tamil for compensation, then the Tamil for to get damaged.
\end{quote}

\subsection*{You'll want to know: \ta{மிதி}}
\textbf{\ta{மிதி}} — \emph{miti} — \textquotedblleft{}to pedal; to tread on\textquotedblright{}.

\begin{quote}
\ta{நான்} \ta{மிதிவண்டியை} \ta{வேகமாக} \ta{மிதித்தேன்}.
\end{quote}

\begin{quote}
\emph{nāṉ mitivaṇṭiyai vēkamāka mitittēṉ.}
\end{quote}

\begin{quote}
I pedalled the bicycle fast.
\end{quote}

Past (I …): \ta{மிதித்தேன்}, \emph{mitittēṉ}.

\ta{மிதிவண்டி}, a bicycle, is a vehicle you tread on.

\begin{cousinweb}
\textbf{\ta{மிதி}} (\emph{miti}) is native Dravidian (DEDR 4861), \textquotedblleft{}to tread on, to trample\textquotedblright{}.
\end{cousinweb}

\begin{grammarlens}[title={What you've built}]
One more word for getting around.
\end{grammarlens}

\subsection*{Your turn}
\begin{itemize}
\raggedright
  \item \emph{Say it:} \emph{miti}
  \item \emph{Say it:} \emph{miti}, once more
  \item \emph{Say it:} \emph{nāṉ mitivaṇṭiyai vēkamāka mitittēṉ.} — I pedalled the bicycle fast.
  \item \emph{Recall:} say the Tamil for compensation, then the Tamil for to get damaged, then say \emph{miti} again
\end{itemize}

\begin{tcolorbox}[breakable,colback=teal!4,colframe=teal!35!black,title={Before you move on}]
What does \ta{மிதி} mean? (\textquotedblleft{}To pedal; to tread on\textquotedblright{}.) Say it once more.
\end{tcolorbox}
\section[\texorpdfstring{\ta{சரக்கு} (carakku)}{சரக்கு (carakku)}]{\ta{சரக்கு} (carakku) — goods; cargo}

\label{lesson:TA-C452-carakku}



\begin{quote}
Before the new one: say the Tamil for to get damaged, then the Tamil for to pedal.
\end{quote}

\subsection*{You'll want to know: \ta{சரக்கு}}
\textbf{\ta{சரக்கு}} — \emph{carakku} — \textquotedblleft{}goods; cargo\textquotedblright{}.

\begin{quote}
\ta{சரக்கு} \ta{ரயில்} \ta{மெதுவாகப்} \ta{போனது}.
\end{quote}

\begin{quote}
\emph{carakku rayil metuvākap pōṉatu.}
\end{quote}

\begin{quote}
The goods train went slowly.
\end{quote}

\ta{சரக்கு} \ta{ரயில்} is a goods train; \ta{சரக்கு} \ta{வண்டி}, a goods vehicle.

\begin{cousinweb}
\textbf{\ta{சரக்கு}} (\emph{carakku}) is native Dravidian (DEDR 2353), \textquotedblleft{}goods, articles of merchandise\textquotedblright{}.
\end{cousinweb}

\begin{grammarlens}[title={What you've built}]
One more word for getting around.
\end{grammarlens}

\subsection*{Your turn}
\begin{itemize}
\raggedright
  \item \emph{Say it:} \emph{carakku}
  \item \emph{Say it:} \emph{carakku}, once more
  \item \emph{Say it:} \emph{carakku rayil metuvākap pōṉatu.} — The goods train went slowly.
  \item \emph{Recall:} say the Tamil for to get damaged, then the Tamil for to pedal, then say \emph{carakku} again
\end{itemize}

\begin{tcolorbox}[breakable,colback=teal!4,colframe=teal!35!black,title={Before you move on}]
What does \ta{சரக்கு} mean? (\textquotedblleft{}Goods; cargo\textquotedblright{}.) Say it once more.
\end{tcolorbox}
\section[\texorpdfstring{\ta{பிரதேசம்} (piratēcam)}{பிரதேசம் (piratēcam)}]{\ta{பிரதேசம்} (piratēcam) — a region; a territory}

\label{lesson:TA-C452-piratecam}



\begin{quote}
Before the new one: say the Tamil for to pedal, then the Tamil for goods.
\end{quote}

\subsection*{You'll want to know: \ta{பிரதேசம்}}
\textbf{\ta{பிரதேசம்}} — \emph{piratēcam} — \textquotedblleft{}a region; a territory\textquotedblright{}.

\begin{quote}
\ta{இந்தப்} \ta{பிரதேசத்தில்} \ta{மழை} \ta{குறைவு}.
\end{quote}

\begin{quote}
\emph{intap piratēcattil maḻai kuṟaivu.}
\end{quote}

\begin{quote}
There is little rain in this region.
\end{quote}

A -\ta{அம்} noun: \ta{பிரதேசத்தில்}, in the region.

\begin{cousinweb}
\textbf{\ta{பிரதேசம்}} (\emph{piratēcam}) is borrowed from Sanskrit \emph{pradeśa}, \textquotedblleft{}a region, a country\textquotedblright{} (CDIAL 8668).
\end{cousinweb}

\begin{grammarlens}[title={What you've built}]
One more word for getting around.
\end{grammarlens}

\subsection*{Your turn}
\begin{itemize}
\raggedright
  \item \emph{Say it:} \emph{piratēcam}
  \item \emph{Say it:} \emph{piratēcam}, once more
  \item \emph{Say it:} \emph{intap piratēcattil maḻai kuṟaivu.} — There is little rain in this region.
  \item \emph{Recall:} say the Tamil for to pedal, then the Tamil for goods, then say \emph{piratēcam} again
\end{itemize}

\begin{tcolorbox}[breakable,colback=teal!4,colframe=teal!35!black,title={Before you move on}]
What does \ta{பிரதேசம்} mean? (\textquotedblleft{}A region; a territory\textquotedblright{}.) Say it once more.
\end{tcolorbox}
\section[\texorpdfstring{\ta{பகுதி} (pakuti)}{பகுதி (pakuti)}]{\ta{பகுதி} (pakuti) — an area; a part}

\label{lesson:TA-C452-pakuti}



\begin{quote}
Before the new one: say the Tamil for goods, then the Tamil for a region.
\end{quote}

\subsection*{You'll want to know: \ta{பகுதி}}
\textbf{\ta{பகுதி}} — \emph{pakuti} — \textquotedblleft{}an area; a part\textquotedblright{}.

\begin{quote}
\ta{நகரத்தின்} \ta{பழைய} \ta{பகுதியில்} \ta{குறுகிய} \ta{தெருக்கள்} \ta{இருக்கின்றன}.
\end{quote}

\begin{quote}
\emph{nakarattiṉ paḻaiya pakutiyil kuṟukiya terukkaḷ irukkiṉṟaṉa.}
\end{quote}

\begin{quote}
There are narrow streets in the old part of the city.
\end{quote}

Also a part of a book or a whole: \ta{முதல்} \ta{பகுதி}, the first part.

\begin{cousinweb}
\textbf{\ta{பகுதி}} (\emph{pakuti}) is native Dravidian (DEDR 3808), from \ta{பகு}, \textquotedblleft{}to divide, to share\textquotedblright{}: a portion.
\end{cousinweb}

\begin{grammarlens}[title={What you've built}]
One more word for getting around.
\end{grammarlens}

\subsection*{Your turn}
\begin{itemize}
\raggedright
  \item \emph{Say it:} \emph{pakuti}
  \item \emph{Say it:} \emph{pakuti}, once more
  \item \emph{Say it:} \emph{nakarattiṉ paḻaiya pakutiyil kuṟukiya terukkaḷ irukkiṉṟaṉa.} — There are narrow streets in the old part of the city.
  \item \emph{Recall:} say the Tamil for goods, then the Tamil for a region, then say \emph{pakuti} again
\end{itemize}

\begin{tcolorbox}[breakable,colback=teal!4,colframe=teal!35!black,title={Before you move on}]
What does \ta{பகுதி} mean? (\textquotedblleft{}An area; a part\textquotedblright{}.) Say it once more.
\end{tcolorbox}
\section[\texorpdfstring{\ta{உருள்} (uruḷ)}{உருள் (uruḷ)}]{\ta{உருள்} (uruḷ) — to roll along; to roll}

\label{lesson:TA-C452-urul}



\begin{quote}
Before the new one: say the Tamil for a region, then the Tamil for an area.
\end{quote}

\subsection*{You'll want to know: \ta{உருள்}}
\textbf{\ta{உருள்}} — \emph{uruḷ} — \textquotedblleft{}to roll along; to roll\textquotedblright{}.

\begin{quote}
\ta{கல்} \ta{மலையிலிருந்து} \ta{கீழே} \ta{உருண்டது}.
\end{quote}

\begin{quote}
\emph{kal malaiyiliruntu kīḻē uruṇṭatu.}
\end{quote}

\begin{quote}
A stone rolled down the hill.
\end{quote}

Past (it …): \ta{உருண்டது}, \emph{uruṇṭatu}.

\ta{ள்} meets the past -\ta{ந்த}- as \ta{ண்ட}: \ta{உருண்டது}. \ta{உருட்டு} is to roll something.

\begin{cousinweb}
\textbf{\ta{உருள்}} (\emph{uruḷ}) is native Dravidian (DEDR 664), \textquotedblleft{}to roll, to tumble, to revolve\textquotedblright{}.
\end{cousinweb}

\begin{grammarlens}[title={What you've built}]
One more word for getting around.
\end{grammarlens}

\subsection*{Your turn}
\begin{itemize}
\raggedright
  \item \emph{Say it:} \emph{uruḷ}
  \item \emph{Say it:} \emph{uruḷ}, once more
  \item \emph{Say it:} \emph{kal malaiyiliruntu kīḻē uruṇṭatu.} — A stone rolled down the hill.
  \item \emph{Recall:} say the Tamil for a region, then the Tamil for an area, then say \emph{uruḷ} again
\end{itemize}

\begin{tcolorbox}[breakable,colback=teal!4,colframe=teal!35!black,title={Before you move on}]
What does \ta{உருள்} mean? (\textquotedblleft{}To roll along; to roll\textquotedblright{}.) Say it once more.
\end{tcolorbox}
\section[(dialogue)]{Getting around, second pass}

\label{lesson:TA-R452-review-446-452}



\begin{quote}
Say the Tamil for an account and a balance.
\end{quote}

\subsection*{Your turn}
Say these aloud:

\begin{itemize}
\raggedright
  \item an account, a balance, interest on money, a cheque, cash
  \item an advance payment, an instalment, an amount due, the remainder, a sum
  \item to save up, to spend money, to waste, to be left over, income
  \item an auction, a lease, a transaction, an investment, to pawn
  \item the original, a fake, luxury, the cost of living, a banknote
  \item a complex, a kilo, a subscription, compensation, to get damaged
  \item to pedal, goods, a region, an area, to roll along
\end{itemize}

\begin{tcolorbox}[breakable,colback=teal!4,colframe=teal!35!black,title={Before you move on}]
An account? (\textbf{\ta{கணக்கு}}, \emph{kaṇakku}.) A balance? (\textbf{\ta{இருப்பு}}, \emph{iruppu}.) Interest on money? (\textbf{\ta{வட்டி}}, \emph{vaṭṭi}.) A cheque? (\textbf{\ta{காசோலை}}, \emph{kācōlai}.) Cash? (\textbf{\ta{ரொக்கம்}}, \emph{rokkam}.) An advance payment? (\textbf{\ta{முன்பணம்}}, \emph{muṉpaṇam}.) An instalment? (\textbf{\ta{தவணை}}, \emph{tavaṇai}.) An amount due? (\textbf{\ta{பாக்கி}}, \emph{pākki}.) The remainder? (\textbf{\ta{மீதி}}, \emph{mīti}.) A sum? (\textbf{\ta{தொகை}}, \emph{tokai}.) To save up? (\textbf{\ta{சேமி}}, \emph{cēmi}.) To spend money? (\textbf{\ta{செலவழி}}, \emph{celavaḻi}.) To waste? (\textbf{\ta{வீணாக்கு}}, \emph{vīṇākku}.) To be left over? (\textbf{\ta{மிஞ்சு}}, \emph{miñcu}.) Income? (\textbf{\ta{வருமானம்}}, \emph{varumāṉam}.) An auction? (\textbf{\ta{ஏலம்}}, \emph{ēlam}.) A lease? (\textbf{\ta{குத்தகை}}, \emph{kuttakai}.) A transaction? (\textbf{\ta{பரிவர்த்தனை}}, \emph{parivarttaṉai}.) An investment? (\textbf{\ta{முதலீடு}}, \emph{mutalīṭu}.) To pawn? (\textbf{\ta{அடகு} \ta{வை}}, \emph{aṭaku vai}.) The original? (\textbf{\ta{அசல்}}, \emph{acal}.) A fake? (\textbf{\ta{போலி}}, \emph{pōli}.) Luxury? (\textbf{\ta{ஆடம்பரம்}}, \emph{āṭamparam}.) The cost of living? (\textbf{\ta{விலைவாசி}}, \emph{vilaivāci}.) A banknote? (\textbf{\ta{நோட்டு}}, \emph{nōṭṭu}.) A complex? (\textbf{\ta{வளாகம்}}, \emph{vaḷākam}.) A kilo? (\textbf{\ta{கிலோ}}, \emph{kilō}.) A subscription? (\textbf{\ta{சந்தா}}, \emph{cantā}.) Compensation? (\textbf{\ta{இழப்பீடு}}, \emph{iḻappīṭu}.) To get damaged? (\textbf{\ta{சேதமடை}}, \emph{cētamaṭai}.) To pedal? (\textbf{\ta{மிதி}}, \emph{miti}.) Goods? (\textbf{\ta{சரக்கு}}, \emph{carakku}.) A region? (\textbf{\ta{பிரதேசம்}}, \emph{piratēcam}.) An area? (\textbf{\ta{பகுதி}}, \emph{pakuti}.) To roll along? (\textbf{\ta{உருள்}}, \emph{uruḷ}.)
\end{tcolorbox}
\section[(dialogue)]{Heading Home — once more}

\label{lesson:TA-R452-review-final-452}



\begin{quote}
Say the Tamil for to pedal and goods.
\end{quote}

\subsection*{Your turn}
\begin{itemize}
\raggedright
  \item \emph{Say it:} to pedal, goods, a region, an area, to roll along
\end{itemize}

\begin{tcolorbox}[breakable,colback=teal!4,colframe=teal!35!black,title={Before you move on}]
To pedal? (\textbf{\ta{மிதி}}, \emph{miti}.) Goods? (\textbf{\ta{சரக்கு}}, \emph{carakku}.) A region? (\textbf{\ta{பிரதேசம்}}, \emph{piratēcam}.) An area? (\textbf{\ta{பகுதி}}, \emph{pakuti}.) To roll along? (\textbf{\ta{உருள்}}, \emph{uruḷ}.)
\end{tcolorbox}
